%% ============================================================
%%  Chapitre 4 — Moteur WFO des Signaux
%%  (Walk-Forward Optimization par catégorie d'indicateurs)
%% ============================================================
\chapter{Moteur WFO des Signaux — Optimisation Dynamique par Catégorie}
\label{ch:wfo}

%% ============================================================
\section{Motivation — Pourquoi un Moteur WFO Distinct du Signal Engine~?}
\label{sec:wfo_motivation}

Le Signal Engine A→G (Chapitre~\ref{ch:signaux}) évalue la robustesse temporelle
d'une variante par validation sur des sous-périodes \emph{avec des paramètres fixes}.
Cette approche convient pour construire un catalogue de signaux fiables, mais elle
présente une limitation structurelle~: les paramètres optimaux d'un indicateur
ne sont pas stationnaires dans le temps.

Le \textbf{Moteur WFO des Signaux} résout ce problème en appliquant la logique de
Walk-Forward Analysis (WFA) de Pardo~\cite{pardo2008} à la \emph{sélection d'indicateurs}
elle-même. Pour chaque catégorie (Tendance, Momentum, Oscillation, Volume) et chaque
horizon (court, moyen, long), le moteur~:

\begin{enumerate}
    \item Optimise les paramètres d'indicateur sur une fenêtre In-Sample (IS) en
    maximisant une fonction objectif robuste (PROM)~;
    \item Teste les paramètres sélectionnés sur la fenêtre Out-Of-Sample (OOS) suivante~;
    \item Répète ce cycle sur plusieurs fenêtres glissantes~;
    \item À partir de la dernière fenêtre, sélectionne un ensemble décorrélé de
    représentants pondérés par leur PROM lissé~;
    \item Calcule le signal courant et les métriques de robustesse WFO.
\end{enumerate}

Ce moteur est \textbf{distinct} du moteur de Backtest WFO classique (qui, lui, optimise
une \emph{stratégie complète} avec gestion des risques et tailles de position). Le
Moteur WFO des Signaux opère exclusivement à l'étape de sélection d'indicateurs
(Station~2).

%% ============================================================
\section{Architecture du Pipeline WFO des Signaux}
\label{sec:wfo_architecture}

\figuretodo{ch4_wfo_pipeline}{Pipeline WFO des Signaux pour une catégorie donnée :\\\
Axe temps (2016 → 2025). Fenêtres IS/OOS glissantes annotées.\\\
Étape 1 : Grille de variants → Calcul PROM IS → Sélection du gagnant IS.\\\
Étape 2 : Application OOS → Calcul oos\_return, oos\_sharpe.\\\
Répétition sur $N$ fenêtres. Dernière fenêtre : sélection ensemble décorrélé.\\\
Sortie : score\_pct $\in [-100, +100]$, robustness\_grade, composite\_score.}

\subsection{Grille de candidats par catégorie}

La grille de candidats est construite en unissant toutes les familles de la
catégorie sur les 30~valeurs de paramètres définies en Layer~A
(Chapitre~\ref{ch:signaux}). Pour la catégorie \texttt{tendance}~:

\begin{equation}
\text{Pool}_{\text{tendance}} = \text{SMA} \cup \text{EMA} \cup \text{EMA-Cross}
\cup \text{Ichimoku} \cup \text{PSAR}
\label{eq:pool_tendance}
\end{equation}

Chaque famille contribue typiquement 30 variantes par horizon, donnant un pool de
100~à~200 candidats selon la catégorie.

\subsection{Politique de fenêtres stricte (\textit{strict\_fold\_driven})}

Le moteur utilise une politique de sélection automatique des fenêtres conçue pour
maximiser le nombre de folds OOS tout en respectant les contraintes de l'horizon~:

\begin{itemize}
    \item \textbf{Bande de train}~: pour chaque horizon, l'étendue des fenêtres
    d'entraînement est bornée~: court [168, 336] barres, moyen [378, 630] barres,
    long [672, 1\,008] barres~;
    \item \textbf{Ratio IS/OOS}~: 25~à~35~\% (ancrages 0.25, 0.30, 0.35)~;
    \item \textbf{Minimum de folds}~: au moins 5 fenêtres OOS sont requises par défaut~;
    \item \textbf{Fallback progressif}~: si le minimum n'est pas atteignable, le seuil
    est abaissé progressivement jusqu'à un plancher de 3 folds.
\end{itemize}

Parmi toutes les configurations candidates, la meilleure est sélectionnée selon~:

\begin{equation}
\text{best} = \arg\min \Bigl(-\text{WFE},\; -\text{robustness\_ratio},\;
\text{dominance},\; \text{distance\_anchor},\; -\text{total\_oos\_pnl}\Bigr)
\label{eq:config_ranking}
\end{equation}

%% ============================================================
\section{Évaluation Par Fenêtre~: La Fonction Objectif PROM}
\label{sec:prom}

\subsection{Origine et Principe}

Le \textbf{PROM} (Pessimistic Return on Margin) est la fonction objectif de
rédérence de Pardo~\cite{pardo2008} pour comparer des stratégies de trading lors
d'une Walk-Forward Analysis. L'idée centrale est la suivante~:

\begin{quote}
\textit{Le profit net brut récompense autant une stratégie qui réalise
100 petits trades rentables qu'une stratégie qui réalise un seul trade
exceptionnel. Le PROM, lui, les distingue~: il pénalise la seconde car
sa performance est non-répétable et statistiquement fragile.}
\end{quote}

Le PROM est dit «~pessimiste~» car il applique une pénalité croissante au
fur et à mesure que le nombre de trades gagnants (resp. perdants) est faible~:
avec peu de trades, la performance observée peut être due au hasard, et le
PROM en réduit proactivement l'estimation.

\subsection{Formule Implémentée}

Pour une liste de retours de trades $\{r_1, \ldots, r_N\}$, on sépare
gains et pertes~:
\begin{align}
\mathcal{W} &= \{r_i \mid r_i > 0\}, \quad N_W = |\mathcal{W}| \label{eq:prom_wins}\\
\mathcal{L} &= \{r_i \mid r_i \leq 0\}, \quad N_L = |\mathcal{L}| \label{eq:prom_losses}
\end{align}

Les gains ajustés et les pertes ajustées sont calculés avec une pénalité
en racine carrée sur le nombre de trades~:

\begin{align}
\overline{W}_{\text{adj}} &= \bar{w} \cdot \bigl(N_W - \sqrt{N_W}\bigr)
\label{eq:prom_adj_wins}\\
\overline{L}_{\text{adj}} &= \bar{l} \cdot \bigl(N_L + \sqrt{N_L}\bigr)
\label{eq:prom_adj_losses}
\end{align}

où $\bar{w} = \text{moyenne}(\mathcal{W})$ et $\bar{l} = \text{moyenne}(|\mathcal{L}|)$.
Le PROM normalisé par le capital $K$ (ou le gross profit simulé) est~:

\begin{equation}
\boxed{\text{PROM} = \frac{\overline{W}_{\text{adj}} - \overline{L}_{\text{adj}}}{K}}
\label{eq:prom}
\end{equation}

\subsection{Lecture et Interprétation}

La pénalité $\sqrt{N}$ a une signification statistique directe~: elle reflète
l'incertitude sur la moyenne d'un échantillon de taille $N$ (proportionnelle à
$1/\sqrt{N}$). En pratique~:

\begin{itemize}
    \item \textbf{PROM $\gg 0$}~: la variante génère des profits réguliers,
    robustes au nombre de trades. C'est le signal à retenir.
    \item \textbf{PROM $\approx 0$}~: la variante est à l'équilibre~;
    ses rendements ne justifient pas le risque transactionnel.
    \item \textbf{PROM $< 0$}~: la variante perd de l'argent sur la fenêtre IS.
    Elle est écartée lors de la sélection des représentants.
    \item \textbf{PROM élevé avec peu de trades}~: le PROM est
    automatiquement pénalisé~; une stratégie avec $N_W = 4$ trades gagnants
    se voit appliquer un coefficient $4 - \sqrt{4} = 2$ au lieu de~$4$,
    soit une réduction de 50~\% de son gain apparent.
\end{itemize}

\begin{table}[H]
\centering
\caption{Exemple d'effet de la pénalité $\sqrt{N}$ du PROM}
\label{tab:prom_exemple}
\begin{tabularx}{\linewidth}{@{}lYYYY@{}}
\toprule
\textbf{Variante} & $N_W$ & $\bar{w}$ & $\overline{W}_{\text{adj}}$
& \textbf{Interprétation} \\
\midrule
SMA(50) — 40 trades & 30 & 1.2\,\% & $1.2\% \times (30 - 5.5) = 29.4\%$
& Robuste (grand $N$) \\
SMA(200) — 6 trades & 5 & 2.1\,\% & $2.1\% \times (5 - 2.2) = 5.9\%$
& Pénalisé (petit $N$) \\
\bottomrule
\end{tabularx}
\end{table}

Dans l'implémentation, pour chaque variante $v$ et chaque fenêtre IS,
la liste des trades est extraite à partir du tableau de signal journalier
(positions longues/flat), et le PROM est calculé sur les retours trade-par-trade
net des coûts de transaction. Un lissage par voisinage (\textit{neighbor-averaging})
est ensuite appliqué au PROM brut pour réduire la sensibilité aux pics isolés
lors du classement final.

\subsection{Coût de transaction}

Un modèle de coût réaliste est intégré dans l'évaluation~: un coût bilatéral
de 10~points de base (bps) par aller-retour est appliqué par défaut, calibré
sur les conditions du marché marocain. Ce coût est soustrait des rendements
journaliers~:

\begin{equation}
r_t^{\text{net}} = s_{t-1} \cdot \frac{P_t - P_{t-1}}{P_{t-1}} -
\left|s_t - s_{t-1}\right| \cdot \frac{c}{10\,000}
\label{eq:cost_model}
\end{equation}

où $s_t \in \{-1, 0, +1\}$ est le signal au temps $t$ et $c$ le coût en bps.

%% ============================================================
\section{Construction de l'Ensemble Représentant}
\label{sec:wfo_ensemble}

À la fin du cycle de walk-forward, les variantes sont classées par leur PROM
lissé sur la dernière fenêtre IS. La sélection de l'ensemble suit un algorithme
glouton décorrélant~:

\begin{enumerate}
    \item Trier les variantes par PROM lissé décroissant~;
    \item Sélectionner la variante de tête comme premier représentant~;
    \item Pour chaque variante suivante~: calculer la corrélation de Pearson avec
    tous les représentants déjà sélectionnés sur les 252 dernières barres~;
    \item Accepter si $\max_j |r_{ij}| \leq 0.85$, rejeter sinon~;
    \item Arrêter dès que \texttt{max\_reps} représentants sont atteints.
\end{enumerate}

Les représentants sélectionnés sont pondérés proportionnellement à leur PROM lissé
(plancher à $0.01$ pour éviter les poids nuls), et leurs signaux courants sont
calculés sur la dernière barre disponible~:

\begin{equation}
\text{score\_pct} = 100 \times \frac{\displaystyle\sum_{i} s_i \cdot \text{PROM}_i}
{\displaystyle\sum_{i} \text{PROM}_i}
\quad \in [-100,\; +100]
\label{eq:wfo_score}
\end{equation}

%% ============================================================
\section{Métriques de Robustesse et Système de Notation}
\label{sec:wfo_scoring}

\subsection{Les quatre métriques fondamentales}

Pour chaque couple (symbole, horizon, catégorie), le moteur produit quatre métriques
agrégées sur l'ensemble des fenêtres OOS~:

\begin{description}
    \item[\textbf{Walk-Forward Efficiency (WFE)}] — Ratio du rendement OOS annualisé
    sur le rendement IS annualisé. Mesure la capacité du système à généraliser
    hors-échantillon. Un WFE $\geq 50\,\%$ est le seuil minimal de Pardo~\cite{pardo2008}~:

    \begin{equation}
    \text{WFE} = \frac{\bar{r}_{\text{OOS}}}{\bar{r}_{\text{IS}}}
    \label{eq:wfe}
    \end{equation}

    \item[\textbf{Robustness Ratio}] — Fraction des fenêtres OOS avec un rendement
    positif~:

    \begin{equation}
    \text{robustness\_ratio} = \frac{\#\{k : r_k^{\text{OOS}} > 0\}}
    {N_{\text{fenêtres}}}
    \label{eq:robustness_ratio}
    \end{equation}

    \item[\textbf{Sharpe OOS moyen}] — Moyenne des ratios de Sharpe annualisés
    calculés sur chaque fenêtre OOS séparément.

    \item[\textbf{Pire drawdown}] — Minimum des rendements OOS par fenêtre~;
    mesure du risque en scénario adverse.
\end{description}

\subsection{Le Score Composite (0–100)}

Ces quatre métriques sont combinées en un score composite pour le classement
inter-catégories~:

\begin{equation}
\text{Composite} = 100 \times \Bigl(
0.40 \cdot \tilde{\text{WFE}} +
0.30 \cdot \tilde{\text{Rob}} +
0.20 \cdot \tilde{\text{Sharpe}} +
0.10 \cdot \tilde{\text{DD}}
\Bigr)
\label{eq:composite}
\end{equation}

où chaque composante est normalisée dans $[0, 1]$~:

\begin{align}
\tilde{\text{WFE}}    &= \text{clamp}(\text{WFE},\; 0,\; 1) \label{eq:norm_wfe} \\
\tilde{\text{Rob}}    &= \text{clamp}(\text{robustness\_ratio},\; 0,\; 1) \label{eq:norm_rob} \\
\tilde{\text{Sharpe}} &= \text{clamp}(\bar{S}_{\text{OOS}} / 2.0,\; 0,\; 1) \label{eq:norm_sharpe} \\
\tilde{\text{DD}}     &= \text{clamp}(1 - |\text{worst\_dd}|,\; 0,\; 1) \label{eq:norm_dd}
\end{align}

\subsection{Le Grade de Robustesse (A–F)}
\label{sec:wfo_grade}

Le grade de robustesse est la synthèse binaire des deux métriques les plus importantes
(WFE et robustness\_ratio). Il constitue l'indicateur principal affiché dans le
tableau de bord~:

\begin{table}[H]
\centering
\caption{Système de notation WFO — Grade de robustesse par catégorie}
\label{tab:wfo_grades}
\begin{tabularx}{\linewidth}{@{}cYYY@{}}
\toprule
\textbf{Grade} & \textbf{Condition WFE} & \textbf{Condition Robustesse} & \textbf{Interprétation} \\
\midrule
\textbf{A} & $\geq 65\,\%$ & \textbf{ET} $\geq 75\,\%$ & Signal très robuste, déployable \\
\textbf{B} & $\geq 55\,\%$ & \textbf{ET} $\geq 60\,\%$ & Signal robuste, fiable \\
\textbf{C} & $\geq 50\,\%$ & \textbf{ET} $\geq 50\,\%$ & Signal acceptable (seuil Pardo) \\
\textbf{D} & $\geq 40\,\%$ & \textbf{OU} $\geq 40\,\%$ & Signal fragile, vigilance requise \\
\textbf{F} & $< 40\,\%$ & \textbf{ET} $< 40\,\%$ & Signal non viable, à ignorer \\
\bottomrule
\end{tabularx}
\end{table}

\textbf{Règle opérationnelle~:} seules les catégories de grade A, B ou C sont
comptabilisées dans le calcul du consensus global (le champ
\texttt{categories\_viable} du résultat global).

%% ============================================================
\section{Signal Global et Modulation S/R}
\label{sec:wfo_global}

\subsection{Consensus pondéré par le score composite}

Les signaux des quatre catégories sont agrégés en un score global pondéré~:

\begin{equation}
\text{score\_global} = \sum_{\text{cat} \in \text{Succ.}}
\frac{\text{composite\_cat}}{\sum_j \text{composite}_j} \cdot \text{score\_cat}
\label{eq:wfo_global_score}
\end{equation}

Cette pondération par le score composite garantit que les catégories avec un WFE
élevé et une robustesse forte ont plus de poids dans le signal final.

\subsection{Modulation par les niveaux S/R}

Le score brut est ensuite modulé par la proximité du prix aux niveaux de
support et de résistance calculés algorithmiquement (méthode \textit{swing levels}
primaire, pivot points en fallback)~:

\begin{equation}
\text{score\_final} = \text{clamp}\bigl(\text{score\_brut} \times m_{\text{S/R}},\; -100,\; +100\bigr)
\label{eq:sr_modulation}
\end{equation}

où le modificateur $m_{\text{S/R}} \in [0.5,\; 1.5]$ amplifie (+15\,\%) le signal
quand il est confirmé par les niveaux, et l'atténue (--15\,\%) quand il les contredit~:

\begin{itemize}
    \item \textbf{Boost $\times 1.15$}~: signal haussier + prix proche d'un support
    (le support confirme l'achat) ou signal baissier + prix proche d'une résistance~;
    \item \textbf{Atténuation $\times 0.85$}~: signal haussier + prix proche d'une
    résistance (risque de blocage) ou signal baissier + prix proche d'un support.
\end{itemize}

\subsection{Tableau des recommandations}

\begin{table}[H]
\centering
\caption{Classification du signal global WFO en recommandations opérationnelles}
\label{tab:wfo_labels}
\begin{tabularx}{\linewidth}{@{}lYY@{}}
\toprule
\textbf{Plage du score global} & \textbf{Label} & \textbf{Recommandation} \\
\midrule
$> +50$ & Très haussier & \texttt{achat\_fort} \\
$+15$ à $+50$ & Haussier & \texttt{achat} \\
$-15$ à $+15$ & Neutre & \texttt{neutre} \\
$-50$ à $-15$ & Baissier & \texttt{vente} \\
$< -50$ & Très baissier & \texttt{vente\_forte} \\
\bottomrule
\end{tabularx}
\end{table}

Ces seuils sont identiques à ceux du Signal Engine A→G (Layer~G,
équation~\eqref{eq:ensemble}), garantissant une homogénéité d'affichage
dans le tableau de bord.

%% ============================================================
\section{Rafraîchissement Automatique et Architecture Asynchrone}
\label{sec:wfo_refresh}

Le Moteur WFO des Signaux est computationnellement intensif~: un calcul complet
(4 catégories × 3 horizons) sur l'ensemble des 93 tickers prend plusieurs heures.
L'architecture asynchrone distingue deux modes~:

\begin{description}
    \item[\textbf{Calcul complet} (\texttt{run\_wfo\_for\_symbol\_horizon})]
    Reconstruit l'ensemble du pipeline WFO pour toutes les catégories
    d'un symbole × horizon. Déclenché une fois par semaine via le batch
    \textit{hebdomadaire} (stratégie \textit{weekly-stale}), ou manuellement
    depuis l'interface.

    \item[\textbf{Rafraîchissement léger} (\texttt{refresh\_wfo\_for\_symbol\_horizon})]
    Re-calcule uniquement les signaux courants des représentants déjà persistés
    (sans re-optimiser les fenêtres WFO). Déclenché automatiquement après chaque
    mise à jour des données de marché (18h00, après clôture). Ce mode permet de
    maintenir les signaux à jour quotidiennement à coût nul.
\end{description}

Les résultats sont persistés dans la base de données (\texttt{WfoSignalSummary}
par catégorie, \texttt{WfoGlobalSignal} pour le consensus) et accessibles via
l'API FastAPI, permettant au tableau de bord de les afficher sans recalcul.

%% ============================================================
\section{Relation avec le Signal Engine A→G}
\label{sec:wfo_vs_se}

Les deux moteurs sont \textbf{complémentaires}, non concurrents~:

\begin{table}[H]
\centering
\caption{Comparaison Signal Engine A→G vs Moteur WFO des Signaux}
\label{tab:se_vs_wfo}
\begin{tabularx}{\linewidth}{@{}lYY@{}}
\toprule
\textbf{Caractéristique} & \textbf{Signal Engine A→G} & \textbf{Moteur WFO des Signaux} \\
\midrule
Optimisation des paramètres & Non (grille fixe) & Oui (PROM IS, fenêtres glissantes) \\
Fenêtres d'évaluation & Sous-périodes fixes & Walk-Forward adaptatif \\
Poids des représentants & Reliability score (Sharpe, stabilité) & PROM lissé \\
Granularité & Par famille d'indicateurs & Par catégorie (4 familles agrégées) \\
Fréquence de calcul & À la demande & Hebdomadaire (complet) + quotidien (léger) \\
Score de robustesse & Reliability score $\in [0, 1]$ & Grade A–F + Composite $\in [0, 100]$ \\
Signal global & Score par famille → consensus G & Score par catégorie → consensus S/R \\
\bottomrule
\end{tabularx}
\end{table}

Le tableau de bord propose les deux vues en parallèle~: l'onglet \textbf{Signal Engine}
donne la profondeur analytique (drill-down par variante, score de stabilité
multi-composante), tandis que l'onglet \textbf{WFO} donne la perspective de
robustesse WFA (grade A–F, ratio OOS/IS, modulation S/R).

\section*{Conclusion du chapitre}

Ce chapitre a présenté le Moteur WFO des Signaux, un mécanisme qui applique la
rigueur de la Walk-Forward Analysis de Pardo non pas à une stratégie complète,
mais à la \emph{sélection dynamique d'indicateurs} par catégorie. Ses contributions
distinctives sont le score composite multi-métriques (WFE, robustness\_ratio,
Sharpe OOS, drawdown), le système de notation A–F directement actionnable, la
modulation par les niveaux S/R, et l'architecture asynchrone permettant une mise
à jour quotidienne sans recalcul complet. Le chapitre suivant présente l'architecture
technique qui héberge et orchestre ces deux moteurs.
